% SPDX-License-Identifier: Apache-2.0
\section{A Unit of Units, Lowered Whole}
\label{sec:x-compose}

A unit whose fields are units. \code{Top} holds a station of three
inputs and the stream ALU of the section before; its \code{run}
makes the channel between them with \code{chan} and joins the two
children's \code{run}s, handing each its ports, exactly as the
hierarchy example does for the simulation. \code{\#[lower]} reads
that \code{run} too: a \code{run} with no loop is a unit of units,
its \code{let}s of \code{chan} and \code{signal} are the nets, and
each \code{self.child.run(..)} is an instance of the child's own
module, its ports joined in order to the nets and to the parent's
ports. The netlist is one module holding two, and the three are one
file in either language; the top is checked against this run under
both simulators at its ports.

The station gathers each request from three channels: the operations
come in request order, the first operands from the last request back,
the second operands odd tags first, with a rest every third cycle on
one of them. A line fires when its three cells hold, and the ALU
answers in the order its units finish, so the results come out in
neither the order the requests were numbered nor the order the
station fired them. The sink holds off for three cycles, and both
children hold with it.

\begin{figure}[htbp]
\centering
\input{compose_timing.tex}
\caption{The composed unit: three inputs in three orders, the
station's masks, the ALU's slots, and the results out of order.}
\label{fig:x-compose}
\end{figure}

\lstinputlisting[caption={\code{ex\_compose.rs}. Run with \code{bazel run //lib/examples:ex\_compose}.},label={lst:x-compose}]{lib/examples/ex_compose.rs}

\lstinputlisting[style=ascii,caption={What \code{ex\_compose} printed when the build ran it: what each source sent, the station's masks, the ALU's slots and the result, per cycle; then the lowering, three modules.},label={lst:x-compose-out}]{out_compose.txt}
